% Scientific Reports Article Methods (placed after the Discussion, as the
% journal suggests; no word limit). Derived from 05_methods_CM.tex:
%  (1) the FuzzyArcLoss V2 companion manuscript (under review, no preprint)
%      is mentioned in the text only, since the journal does not allow
%      unpublished work in the reference list; the published origin of the
%      loss family is cited instead;
%  (2) the tab_atom_operationalization float is in the Supplementary
%      Information (it is adjudication apparatus, not a headline result);
%  (3) the ethics statement and the statement on the use of large language
%      models are in this section, as the journal requires;
%  (4) the six-level interpretability stack and the LLM explanation layer
%      are condensed into one "Released per-slide tooling" paragraph: they
%      are released software, not an evaluated result of this study.
\section*{Methods}

% =====================================================================
\subsection*{Cohort and labels}
% =====================================================================

\paragraph{TCGA-LUAD acquisition.}
We obtained $1{,}054$ haematoxylin-and-eosin (H\&E) WSIs
from the LUAD collection of The Cancer Genome Atlas via the
GDC portal.\cite{tcga2014luad}
After quality filtering (Otsu tissue mask, pen-marker detector,
blur threshold; see \emph{Slide quality control} below), $687$
analysable images remained. That figure is an inventory rather than a
cohort, and three exclusions separate the two. The GDC LUAD project
directory contains slides from patients whose MAF assigns them to
LUSC; cross-referencing barcodes against the LUAD and LUSC mutation
tables removes $150$ slides from squamous patients and $17$ from
patients with no mutation call for any of the six target genes.
Sample-type code \texttt{11A} in the TCGA barcode marks
solid-tissue-normal sections, which carry the patient's germline
barcode and therefore inherit a mutation label while containing no
tumour; $62$ such slides were removed. What remains is $458$ tumour
slides from $444$ patients across $33$ tissue-source sites, and this
is the cohort behind every LUAD number we report. Nine patients
contributed more than one slide, which accounts for the excess of $14$
slides over patients; cross-validation is performed at the
\emph{patient} level, and repeated under site-preserved folds, to
prevent leakage through either unit.

The retained slides are not restricted to diagnostic FFPE blocks:
$183$ are FFPE (DX) and $275$ are frozen (TS/BS) sections, so the
cohort is $40\%$ FFPE. We retain the mixture deliberately and state it
plainly, because a frozen majority is a transferability hazard rather
than a neutral composition detail. Three considerations decide it.
First, arithmetic: filtering to FFPE leaves $183$ slides, at which
\textit{RBM10} would carry roughly $13$ positives and four of six genes
would fall below the estimation floor defined below, so the
FFPE-only cohort cannot support per-gene adjudication at all and the
choice is between a mixed cohort and no cohort.
Second, direction of bias: frozen sections carry freezing artefact,
poorer nuclear detail and section-thickness variability, all of which
degrade the growth-pattern vocabulary more than the unconstrained
encoder, so the mixture works \emph{against} the pattern channel and
therefore against the complementarity findings rather than for them; a
regime-(i) saturation call is the reading that a frozen majority most
easily manufactures, which is why the site-preserved protocol matters
for TP53 and EGFR specifically.
Third, external validity: CPTAC-LUAD is FFPE throughout, so the
cohort-transfer arm below is itself a preparation-shift test, and the
frozen fraction here is one of the candidate mechanisms for the
collapse it reports, alongside scanner and site signature. We do not
claim to separate preparation from site, and say so where the
transfer is discussed.

\paragraph{Target genes.}
Six driver mutations were retained on the basis of
LUAD-relevance and prevalence $\geq 5\%$ in the cohort
(\cref{tab:mutation_prevalence}), quoted over the $458$-slide cohort:
\textit{TP53} (49.6\%, $n_+\!=\!227$), \textit{EGFR}
(15.5\%, $71$), \textit{KRAS} (27.5\%, $126$), \textit{STK11}
(12.9\%, $59$), \textit{KEAP1} (18.6\%, $85$) and
\textit{RBM10} (7.0\%, $32$). The lower rates that appear against the
$687$-image inventory are diluted by the squamous cases and normal
sections excluded above.

The corresponding mutation calls (MAF files) were obtained from
GDC and only non-silent variants retained: missense, nonsense,
frameshift, splice-site, in-frame insertion/deletion.
Synonymous and intronic variants were excluded.

\paragraph{Slide quality control.}
Each WSI was checked for: (i)~presence of a tissue area
$\geq 20\,\text{mm}^{2}$; (ii)~absence of large pen-marker
regions ($>10\%$ of the tissue area, detected by an HSV
band-pass); (iii)~per-tile blur (variance of the Laplacian)
above an empirically fixed threshold; (iv)~absence of large
histology artefacts (tears, folds, air bubbles).
Slides failing any criterion were excluded.

% =====================================================================
\subsection*{Tile extraction and tissue mask}
% =====================================================================

Each WSI was tiled at level~0 into non-overlapping
$224 \times 224$~px patches, that is, at each slide's native scanning
resolution, with no resampling to a common microns-per-pixel.
The analysable cohort is not uniform in that respect: of the $458$
slides, $76.4\%$ were scanned at $40\times$
($\mathrm{mpp}\!<\!0.35$) and $22.6\%$ at $20\times$
($0.35\!\leq\!\mathrm{mpp}\!<\!0.75$; the remaining $1\%$ carry no
resolution metadata), so a fixed $224$~px window spans
about $56~\mu\text{m}$ of tissue on the former and about
$112~\mu\text{m}$ on the latter, and \ctranspath{}---pretrained at
$20\times$\cite{wang2022ctranspath}---sees the majority of our tiles at
twice its training linear scale.
Training and cross-validation run over this same mixture, so the
within-cohort estimates reported here are internally consistent; the
consequence is for transfer to a cohort with a different magnification
profile, which we measure directly in the field-of-view control of
\cref{tab:cptac_refit} rather than assume.
A per-slide Otsu tissue mask discarded blank or background tiles
($<10\%$ tissue).
Per-tile pen-marker, blur and artefact checks were re-applied at
the tile level.
A 4-channel input ($\text{RGB}\,+\,\text{tissue-mask}$) was passed
to the encoder.\cite{wang2022ctranspath}

% =====================================================================
\subsection*{Pattern classifier (Artefact 1)}
% =====================================================================

The six-pattern classifier was trained on the Zenodo-ANORAK
dataset\cite{anorak2024,ANORAKZenodo}---$732$ expert-annotated
tiles spanning lepidic, acinar, papillary, micropapillary,
solid and cribriform, reduced to $637$ after exclusion of
``other tissue'' (class~0) tiles---using a \ctranspath{}
backbone\cite{wang2022ctranspath} and a fuzzy angular-margin head
(FuzzyArcLoss~V2, the computational-pathology extension of the
angular-margin loss family introduced in
ref.~\cite{limareina2025fuzzyarc}; its full derivation and loss-family
ablation are reported in a companion manuscript under review, which is
cited here in the text only, as the journal's reference policy
requires for unpublished work).
Per-sample fuzzy modulation of the ArcFace
margin\cite{deng2019arcface} is calibrated by Optuna
($s^{*}=46.1$, $m^{*}=0.518$, $\tau^{*}=0.448$,
$50$~trials, TPE sampler, MedianPruner).
Training protocol: AdamW (lr~$1\times10^{-4}$), cosine schedule
to $10^{-6}$, $5$-epoch warm-up, $120$ epochs, batch size $32$,
fp32; MixUp disabled.
Reported macro-F1 = $92.31\% \pm 2.04$ across 15 runs (5-fold $\times$
3 seeds) with $p = 0.001$ vs SphereFace and Cohen's $d = 1.00$.
Tile probabilities $\mathbf{p}\in\Delta^5$ become the pattern
channel used downstream.

% =====================================================================
\subsection*{MIL aggregation: PI-ABMIL and FC-MIL}
% =====================================================================

\paragraph{Tile representation.}
For each slide we retained the top $200$ tiles by encoder
embedding norm to keep the bag size constant.
Each tile contributed a $512$-d \ctranspath{} embedding
$\mathbf{e}\in\mathbb{R}^{512}$ and a $6$-d pattern probability
$\mathbf{p}\in\Delta^5$.
Four conditions were ablated as inputs:
B2 ($512$-d embeddings only to
gated ABMIL\cite{ilse2018attention}); B3 ($6$-d pattern
probabilities only); PI-ABMIL ($512+6$ concat); and
FC-MIL ($512$-d embeddings to gated
attention $+$ $6$-d patterns to a fuzzy Choquet head).
Two further conditions of the superseded protocol---B1, a $7$-d
slide-level pattern-fraction and entropy vector to
XGBoost, and A, a one-hot of
$\arg\max\mathbf{p}$---were not scored in the nested-cross-validation
run and are reported nowhere in this paper. They are named here only so
that the four-condition grid is not mistaken for the whole of the
original design.

\paragraph{PI-ABMIL.}
Gated attention\cite{ilse2018attention} computes
$\alpha_i = \softmax(\mathbf{v}^\top \tanh(\mathbf{W}_a
[\mathbf{e}_i;\mathbf{p}_i]) \odot \sigmoid(\mathbf{W}_b
[\mathbf{e}_i;\mathbf{p}_i]))$, and the slide representation is
$\mathbf{z} = \sum_i \alpha_i [\mathbf{e}_i;\mathbf{p}_i]$
followed by a linear gene head.
Hidden dim $128$, two-layer attention; $\approx 199{,}553$
parameters total.

\paragraph{FC-MIL (2-additive fuzzy Choquet).}
The slide-level pattern vector $\mathbf{P}\in\Delta^5$ is
aggregated by a differentiable
$2$-additive Choquet
integral~\cite{grabisch1996application}.
A general fuzzy measure over $n\!=\!6$ patterns would require
$2^{n}\!-\!2\!=\!62$ parameters; the 2-additive approximation
reduces this to $6 + \binom{6}{2}\!=\!21$ parameters, fully
determined by singletons and pairs:
\begin{equation}
g(A) = \sum_{k\in A}\phi_k \;+\; \sum_{\{j,k\}\subseteq A} I_{jk},
\qquad A \subseteq \{1,\ldots,6\},
\label{eq:fuzzy_measure}
\end{equation}
where the singletons $\phi_k$ are Shapley
values~\cite{shapley1953value} (intrinsic contribution of pattern
$k$) and $I_{jk}$ are pairwise interaction indices
(positive $\Rightarrow$ synergy, negative $\Rightarrow$ redundancy).
We parameterise the singletons via a softmax
$\phi_k = e^{\psi_k}/\sum_l e^{\psi_l}$ on a learnable vector
$\psi\in\mathbb{R}^{6}$; the interactions are stored as a learnable
upper-triangular matrix with an $L_1$ sparsity prior
$\mathcal{R}_\text{int}=\lambda_I\sum_{j<k}|I_{jk}|$ and a soft
monotonicity penalty
$\mathcal{R}_\text{mono} =
\sum_{A\subseteq B}\max(0,\,g(A)\!-\!g(B))$
sampled by Monte Carlo at every batch.

For a slide-level fuzzy pattern probability vector
$\mathbf{P}\in\Delta^{5}$ sorted in ascending order
$P_{(1)}\!\leq\!P_{(2)}\!\leq\!\cdots\!\leq\!P_{(6)}$, the discrete
Choquet integral is
\begin{equation}
C_g(\mathbf{P}) \;=\; \sum_{k=1}^{6}
\bigl[P_{(k)}-P_{(k-1)}\bigr]\, g\!\left(A_{(k)}\right),
\quad A_{(k)} \!=\! \{(k),\ldots,(6)\},
\label{eq:choquet_integral}
\end{equation}
with $P_{(0)}\!:=\!0$.
The output $\mathbf{c}_s\in\mathbb{R}^{6}$ is concatenated with the
gated-attention embedding $\mathbf{z}_\mathrm{attn}\in\mathbb{R}^{256}$
and projected to the per-gene prediction.
\fcmil{} is the architectural primitive that exposes
pattern--pattern interactions as the $15$ scalars $I_{jk}$ for
direct clinical inspection at inference time.

% =====================================================================
\subsection*{Partial information decomposition of the pattern and
embedding channels}
% =====================================================================

\paragraph{Decomposition.}
The three-regime taxonomy is stated as a claim about how the
information the slide carries about the mutation label is
distributed between the two channels the aggregator receives.
Let $Y\in\{0,1\}$ be the per-gene mutation label, and let the two
sources be the slide-level pattern channel $\mathbf{P}\in\Delta^{5}$
and the \ctranspath{} embedding channel
$\mathbf{E}\in\mathbb{R}^{512}$.
Partial information decomposition
(\pid{})\cite{williams2010pid,bertschinger2014unique} splits the
joint mutual information\cite{cover2006elements} into four
non-negative atoms,
\begin{equation}
I\!\left(\mathbf{P},\mathbf{E};Y\right) \;=\;
\Ired \;+\; \Iuni{\mathbf{P}} \;+\; \Iuni{\mathbf{E}} \;+\; \Isyn ,
\label{eq:pid}
\end{equation}
where $\Ired$ is the information available \emph{redundantly} from
either channel, $\Iuni{\mathbf{P}}$ and $\Iuni{\mathbf{E}}$ the
information available from one channel only, and $\Isyn$ the
information available only from their \emph{joint} state.
What we take from \cref{eq:pid} is its \emph{exhaustiveness}, not its
values. Two sources distribute their information about a target in these
four ways and in no other, so a taxonomy whose categories are the four
atoms is complete by construction: it cannot harbour an unnoticed fifth
regime, and it can name a category before any observed case occupies it.
That is a property of the partition, and it survives the
non-uniqueness of the decomposition---the competing definitions disagree
on how to apportion information among the atoms, not on how many atoms
there are. Where naming one requires a definitional convention we follow
Bertschinger~et~al.\cite{bertschinger2014unique} No numerical atom value
is reported or used to support any claim in this paper, so no result
depends on that choice.
\Cref{eq:pid} makes the three regimes mutually exclusive
predictions rather than post-hoc descriptions: regime~(i) is
$\Ired$-dominated, so any fusion of $\mathbf{P}$ into $\mathbf{E}$
is a no-op; regime~(ii)---\emph{complementarity}---is the
regime in which the pattern channel carries information the encoder
does not, in either of two sub-forms, \emph{additive
complementarity} (ii-a), $\Iuni{\mathbf{P}}>0$ with
$\Isyn\!\approx\!0$, or \emph{non-additive synergy}, equivalently
inter-pattern co-occurrence, (ii-b), $\Isyn>0$, which only a
non-additive aggregator can realise;
and regime~(iii)---\emph{out-of-ontology}---requires only that the
vocabulary-side atoms vanish,
$\Iuni{\mathbf{P}}\!\approx\!\Isyn\!\approx\!0$, and says nothing about
$\Ired$ or $\Iuni{\mathbf{E}}$, either of which may be large. Its strict
sub-state, in which $I(\mathbf{P},\mathbf{E};Y)$ itself sits near the
estimation floor, is a claim of a different order, and we name it
separately as a \emph{modality ceiling}.

\paragraph{The ceiling is the modality, not the encoder.}
Both channels are deterministic functions of the same whole-slide image
$W$: $\mathbf{E}=f(W)$ under the frozen \ctranspath{} encoder and
$\mathbf{P}=g(W)$ under the frozen ANORAK pattern classifier. The data
processing inequality\cite{cover2006elements} therefore bounds the
entire left-hand side of \cref{eq:pid} by a quantity that belongs to the
tissue rather than to any model,
\begin{equation}
I\!\left(\mathbf{P},\mathbf{E};Y\right) \;\leq\; I\!\left(W;Y\right),
\label{eq:modality_ceiling}
\end{equation}
and this is what separates a claim about the \emph{vocabulary} from a
claim about the \emph{modality}. Regime~(i) says $\mathbf{E}$
already attains the bound in \cref{eq:modality_ceiling}, so a second
channel has nothing left to recover; regime~(ii) says $\mathbf{E}$ falls
short of it and $\mathbf{P}$ restores part of the shortfall;
regime~(iii) says $\mathbf{P}$ restores none of it, which condemns the
vocabulary and not the tissue. \textit{STK11} is the illustration: the
encoder alone reaches $0.591$ on the same slides against $0.499$ for the
vocabulary alone, a gap whose interval excludes zero, so the modality
plainly carries something, and the remedy---tumour-infiltrating
lymphocytes, budding, tumour--stroma ratio---is a new vocabulary over
the \emph{same} H\&E rather than a new modality.

Only the strict sub-state is a claim about the modality. When
$I(\mathbf{P},\mathbf{E};Y)$ itself sits near the estimation floor the
bound in \cref{eq:modality_ceiling} is near zero, and then no encoder,
no aggregator and no label budget can move it, because the deficit is in
the tissue rather than in the pipeline. Even there the claim is
evidential rather than proved: what any experiment observes is the
left-hand side of \cref{eq:modality_ceiling}, and a total near the floor
licenses the statement about $I(W;Y)$ only insofar as neither frozen
channel is itself discarding the relevant signal. We therefore report a
modality ceiling as the strongest evidence currently obtainable that one
has been reached rather than as a proof, and its remedy---a source
outside the modality altogether---follows from the evidence either way.
Glioma MGMT is the only target in this work whose total sits at the
floor under every condition. For \textit{KEAP1} the case that the limit
is the modality rather than the vocabulary rests on the biochemical
nature of the NRF2 axis and on both single-channel conditions sitting at
chance, not on a total at the floor, and we present it on that basis.

The same inequality has been invoked for \emph{molecular translators}
in which the translated representation \emph{replaces} its
source.\cite{saisan2026ceilings} Our setting is the augmenting one:
$\mathbf{P}$ is supplied alongside $\mathbf{E}$, so
$I(\mathbf{P},\mathbf{E};Y)$ is bounded by \cref{eq:modality_ceiling}
rather than by $I(\mathbf{E};Y)$, and the gap between the two is the
room in which regime~(ii) lives.

\paragraph{Correspondence with the Choquet interaction index.}
For a $2$-additive fuzzy measure, \cref{eq:fuzzy_measure} assigns
the pair $\{j,k\}$ a value that exceeds the sum of its singletons
exactly when $I_{jk}>0$; the aggregator therefore pays for the
\emph{joint} activation of patterns $j$ and $k$ over and above
their separate
activations.\cite{grabisch1996application}
This is the decision-theoretic counterpart of a positive
synergistic atom in \cref{eq:pid}, and, symmetrically, $I_{jk}<0$
is the counterpart of redundancy between the two patterns.
The sign convention is the axiomatic one: the interaction index of a
pair is positive when the pair is complementary, negative when it is
redundant, and zero under additive independence.
It matters what kind of object $I_{jk}$ is. It is not an estimator of
$\Isyn$ and is never used as one. It is a parameter of our own
aggregator, learned by gradient descent and read off the fitted fuzzy
measure, whose semantics are fixed by the axioms above rather than
inferred from the data. What it records is whether the aggregator found
it worthwhile to pay for the joint activation of a pattern pair over and
above their separate activations---a fact about the fitted decision
function, available exactly. Its bearing on \cref{eq:pid} is
qualitative: an aggregator profits from a conjunction its additive
counterpart cannot express only when there is something in the joint
state to profit from, which is the content of a non-empty synergistic
atom. Symmetrically, the additive concat fusion of \piabmil{} is an
aggregator whose hypothesis space forces any such conjunction to be
discarded, which is what makes the \fcmil{}-versus-\piabmil{} contrast
informative.
Two further per-slide read-outs of the fitted model stand in the same
relation to the other atoms: the SHAP embedding/pattern split tracks
$\Ired$ against $\Iuni{\mathbf{E}}$, and the single-channel ablation
reports what removing a channel costs. Each is a property of the fitted
model that a given atom predicts; none is a measurement of that atom.

\emph{Which of these carries evidential weight.} The four training conditions
do; the per-slide read-outs do not. B2, B3, \piabmil{} and \fcmil{} are
complete retrainings under one protocol, so a contrast between them compares
what each hypothesis space can extract from the data. The SHAP split and
the ablation instead
interrogate a single fitted function: the split attributes that
function's output, and the ablation reports what zeroing an input at
inference costs a model that never saw that input at zero. That measures what
the fitted function depends on, not what the channel contains---a redundant
channel can be load-bearing for one fit and still carry nothing the other
channel lacks. The retrained-without-the-channel form of that question is
already in the design, and it is precisely B2 and B3. We therefore rest every
regime claim on the four conditions and read the per-slide read-outs as
illustrations
of the mechanism a claim asserts, never as its warrant. The single place they
are load-bearing is the (i)-versus-(iii) boundary, where the four conditions
collide by construction; we mark that dependence where it arises rather than
spreading it silently across the taxonomy.

\paragraph{Reading the atoms from the model-side signature.}
We do not estimate the atoms of \cref{eq:pid} numerically, and the governing
reason is identification rather than sample size. Competing decompositions
return different atom values on identical
data\cite{williams2010pid,bertschinger2014unique,liang2023quantifying} and no
agreed axiom set singles one out, so a point value for $\Isyn$ is not an
identified quantity and would not become one at infinite $n$. Sample size
compounds this rather than causing it: a direct estimate on the $512$-d
channel is intractable at $n=444$ under any of those definitions
(\cref{fig:pid_estimator}; the estimator we specify but decline to use
is \cref{tab:pid_estimator}). Reading the
\emph{sign} of contrasts between fully retrained models is what those two
facts license---it is the correct response to a non-identified target, not a
fallback forced by a small cohort. Instead each
atom's \emph{presence} is read from the model-side signatures introduced
above---the sign of the learned Choquet interaction $I_{jk}$ for
$\Isyn$, the SHAP embedding/pattern split for $\Ired$ versus
$\Iuni{\mathbf{E}}$, and the single-channel ablations for the
unique atoms---together with the contrast among the four training
conditions B2/B3/PI/FC (\cref{tab:atom_operationalization}). The regime is
therefore assigned from the
\emph{direction} of these signatures, never from a point estimate of an
atom or from the magnitude of any \auroc{}.

\emph{The noise-aware rule.} Each regime is a claim about one deciding
contrast, so it is that contrast whose uncertainty has to carry the
claim. We resample patients with replacement from the pooled out-of-fold
predictions, recompute every condition's \auroc{} on each draw, and take
the $2.5$th and $97.5$th percentiles of the contrast as its interval.
Resampling rather than retraining is what makes the replicate count
affordable; it quantifies sampling noise in the contrast, which is the
quantity the regime read depends on, and not training-time variance.
Complementarity is then asserted only if its deciding contrast excludes
zero---$\mathrm{PI}-\mathrm{B2}$ for the additive sub-form,
$\mathrm{FC}-\mathrm{PI}$ for the non-additive one. Regimes~(i)
and~(iii) make the opposite claim, that \emph{no} contrast is meaningful,
so neither is a residual category: both require the two fusion-versus-encoder
intervals, $\mathrm{PI}-\mathrm{B2}$ and $\mathrm{FC}-\mathrm{B2}$, to fall
inside $\pm0.02$, which is an equivalence test rather than the mere
absence of a significant contrast, since wide intervals would satisfy the
latter vacuously. The other two contrasts we compute carry different claims
and do not gate equivalence: $\mathrm{FC}-\mathrm{PI}$ decides (ii-b), and
$\mathrm{B3}-\mathrm{B2}$ measures what the vocabulary omits rather than what
fusion adds.
The margin is not a convention. The eight published aggregator gains audited
in \cref{tab:literature_power_audit} average $0.0216$, so $\pm0.02$ is the
effect size this literature typically claims for an aggregator change;
declaring equivalence at that margin is declaring that a fusion arm does not
deliver what the field says such a change delivers. Sweeping
$\delta\!\in\!\{0.01,0.02,0.03\}$ leaves every LUAD regime read unchanged; at
$\delta\!=\!0.05$ two of the twelve cells would become equivalent, but that
margin exceeds the largest gain in the entire audited literature ($0.046$), so
it would certify as equivalent a difference larger than any this field has
claimed (\cref{tab:threshold_sensitivity}). What binds at this cohort size is
interval width rather than the margin: the median cell would need roughly
$14{,}000$ patients before equivalence could be certified at $\pm0.02$, which
is the same order as the conditional regime test
(\cref{tab:design_power}). The framework is therefore adjudicable, but not at
this $n$, and we state that as a design property rather than leaving it for
the reader to infer from null results. A comparison meeting neither
condition returns no regime, which is what every LUAD target and the
CPTAC transfer do (\cref{tab:luad_contrasts}). One read is taken from
levels rather than contrasts: when every condition, the encoder alone
included, has an interval whose upper bound lies below the Conclusive
threshold, no reader of the modality in the design reaches the target,
which is the modality-ceiling form of (iii); it needs no contrast and
is the only regime read in this work that a fusion-versus-encoder
equivalence does not gate (glioma MGMT, \cref{tab:crossdomain_merged}).
This is the rule referred
to as noise-aware throughout, and it is what stops a difference smaller
than the fold-to-fold spread from being promoted to a structural claim.
The learned interaction indices $I_{jk}$ are exempt: they are parameters
of the fitted measure rather than differences between \auroc{}s, so their
sign is read directly, against a floor set by the smallest $|I_{jk}|$
the $L_1$ prior leaves non-zero.

\emph{Sample-size calculation.} Every cohort size quoted in this paper
comes from one of two closed forms, both anchored on the hierarchical
bootstrap at $n_0\!=\!444$ patients and both rescaling its standard error
as $n^{-1/2}$, which is the one assumption they make. The first is the
prospective design calculation of \cref{tab:design_power}: the number of
patients at which a stipulated difference $\delta$ in a contrast would be
detected at two-sided $\alpha\!=\!0.05$ with power $1-\beta\!=\!0.80$,
\begin{equation}
n(\delta) \;=\; n_0\left(\frac{(z_{1-\alpha/2}+z_{1-\beta})\,
  \mathrm{SE}_{n_0}}{\delta}\right)^{2}
\;=\; 444\left(\frac{2.80\,\mathrm{SE}_{444}}{\delta}\right)^{2},
\label{eq:design_n}
\end{equation}
with $\mathrm{SE}_{n_0}$ recovered from the bootstrap percentile interval
as $(\mathrm{ci}_{\mathrm{hi}}-\mathrm{ci}_{\mathrm{lo}})/(2\times1.96)$.
For the conditional contrast of \cref{tab:kras_coactivation}
($\mathrm{FC}-\mathrm{PI}$, high minus low co-activation tercile) the
intervals are $[-0.117,+0.072]$ patient-preserved and $[-0.090,+0.141]$
site-preserved, so $\mathrm{SE}_{444}\!=\!0.048$ and $0.059$. At
$\delta\!=\!0.03$, the mid-range of the aggregator gains this literature
reports, \cref{eq:design_n} gives $444\,(2.80\times0.048/0.03)^2 \approx
9{,}000$ patients patient-preserved and $444\,(2.80\times0.059/0.03)^2
\approx 13{,}400$ site-preserved; at $\delta\!=\!0.05$, $3{,}200$ and
$4{,}800$. The KRAS-positive column multiplies by the $27.9\%$
patient-level prevalence. Inverting \cref{eq:design_n} gives the smallest
effect the present cohort resolves, $2.80\times\mathrm{SE}_{444}
\!=\!0.135$ and $0.165$. The second form is the equivalence version
reported per gene in \cref{tab:threshold_sensitivity}c: a contrast with
point estimate $\hat d$ and bootstrap half-width $h_{n_0}$ is certified
inside $\pm\delta$ once $|\hat d| + h_{n_0}\sqrt{n_0/n} \le \delta$, so
\begin{equation}
n_{\mathrm{eq}} \;=\; n_0\left(\frac{h_{n_0}}{\delta-|\hat d|}\right)^{2},
\qquad |\hat d|<\delta ,
\label{eq:equivalence_n}
\end{equation}
evaluated on the binding of the two gating contrasts. For \textit{TP53}
patient-preserved the binding contrast is $\mathrm{FC}-\mathrm{B2}$,
$-0.015$ $[-0.042,+0.010]$, so $h_{444}\!=\!0.026$ and
$n_{\mathrm{eq}}\!=\!444\,(0.026/0.005)^2\approx12{,}600$; where
$|\hat d|\ge\delta$ no cohort certifies equivalence and the cell is marked
\emph{never}.

What these numbers are for needs stating, because it is the point of
reading regimes rather than scores. Neither equation says anything about
the level of a target: a larger cohort does not raise the \auroc{} of
\textit{STK11} or \textit{KEAP1}, since the level is a property of what
the modality carries about the target, not of how precisely it is
measured. What grows with $n$ is the resolution of the \emph{contrast}
between conditions, that is, the ability to say whether fusion adds to
the encoder, adds nothing, or costs. The regime reads exist to separate
the targets whose problem is data from those whose problem is not:
regime~(iii) in its modality-ceiling form is adjudicated from levels and
is already decided at $n\!=\!444$ (glioma MGMT), and no cohort size
changes it; the LUAD regime reads are undecided because their deciding
contrasts are wider than the margin, and \cref{eq:design_n,eq:equivalence_n}
say exactly how much data would decide them. A cohort-size
recommendation is therefore attached to a contrast, never to a target's
\auroc{}.

\emph{The assignment.} The additive sub-form (ii-a) is declared when
$\mathrm{PI}\!>\!\max(\mathrm{B2},\mathrm{B3})$ while $\mathrm{FC}$ does
not exceed $\mathrm{PI}$ and no pair interaction clears the floor.
Regime~(i) holds when no fusion arm adds to either single channel and
$I_{jk}\!\le\!0$, and regime~(iii) when in addition the vocabulary-side
atoms are empty, in the sense set out below.

The non-additive sub-form (ii-b) carries a stricter criterion. An unconditional
$\mathrm{FC}\!>\!\mathrm{PI}$ with $\max_{j<k}I_{jk}>0$ is not
sufficient, because both can hold while the aggregator gain is unrelated
to the pair the index names: $I_{jk}$ is a parameter of the fitted
measure, shared across bags, and an unconditional contrast averages the
cases where the pair co-activates with those where it does not. We
therefore require the contrast to be \emph{conditional} on the
co-activation the index points to. For the pair $(j,k)$ with the largest
positive $I_{jk}$, per-bag co-activation is $\min(f_j,f_k)$ over the hard
vocabulary assignment, averaged to the patient, and the cohort is cut
into terciles of that quantity; the cut uses the source vocabulary only
and no target-side quantity, so it cannot be tuned to the outcome.
Sub-form (ii-b) is declared when the FC${-}$PI contrast in the high
stratum exceeds that in the low stratum by an interval excluding zero
under the hierarchical bootstrap, and under both the patient-preserved
and site-preserved protocols. Where the interval includes zero we report
the sub-form as unsupported together with the smallest difference the
design could have resolved, so a null is distinguishable from an
untested claim (\cref{tab:kras_coactivation,tab:design_power}).

Note that (ii-b) does not
require the additive arm to beat the encoder, and must not. An additive
head losing what a non-additive head keeps is the signature a
synergistic atom predicts, so gating (ii) on \piabmil{} alone would
exclude exactly the case the sub-form exists to name; the criterion is
stated this way for that reason and not because any gene here exhibits
it. KRAS, the one gene whose learned index names a pair
($I(\text{lepidic}\!\times\!\text{solid})\!=\!+0.032$), fails the
criterion on both counts: the conditional contrast is unresolved
(\cref{tab:kras_coactivation}), and the unconditional ordering is
$\mathrm{FC}\!=\!0.506$ below $\mathrm{PI}\!=\!0.516$, both within a
fold's spread of $\mathrm{B2}\!=\!0.517$. The regime axis remains
orthogonal to the Conclusive/Inconclusive verdict
(\cref{eq:conclusive_verdict}), which is the sole magnitude test.

\emph{Where the signature map is not injective.} The map from atom
configuration to model-side signature is not one-to-one, and two collisions
are consequential enough that we name them rather than leave them to be found.

The first is internal to the taxonomy. Regimes~(i) and~(iii) share the
same four-condition signature---no fusion arm adds anything, no pair
interaction survives---and are separated by the single-channel
evidence rather than by the fusion contrasts: under (i) each channel
alone carries the target and the SHAP split places a shared quantity on
both, whereas under (iii) the vocabulary channel alone is at chance and
zeroing it costs the fused model nothing. Because $\Ired$ can only be large where
the target is predictable at all, this one boundary is not independent of
the magnitude axis, and we state that rather than claiming orthogonality
for the taxonomy as a whole: what the KRAS case establishes is that
regime~(ii) is verdict-independent, not that every boundary is. For
\textit{KEAP1} the separation is additionally supported from outside the
model, since the NRF2 axis the gene acts on is biochemical rather than
morphological, which is evidence no fitted model supplies.

The second collision conditions every reading of an empty vocabulary-side
atom, and it is not resolvable from the four conditions at all: a
\emph{redundant} atom seen through an encoder that represents the vocabulary
and an \emph{absent} atom seen through an encoder blind to it produce the same
table. Separating them requires an independent measurement of whether the
encoder represents the vocabulary, and the reading here rests on one. On the
annotated tile benchmark that defines the ANORAK
vocabulary\cite{anorak2024,ANORAKZenodo}, heads trained on the frozen
\ctranspath{} embedding recover the six growth patterns at $90.8$--$92.4\%$
macro-$F_1$ across eleven objectives, with the two most contested classes
separable at one-versus-one $\auroc$ $0.96$--$1.00$; a ViT-H/14 foundation
encoder pretrained by a different group on a different corpus reaches
$90.8$--$92.0\%$ on the same folds under a cached linear
probe (companion manuscript under review, cited above). The embedding space already contains the
vocabulary, and a modern encoder does not contain it better. That is what
licenses reading $\mathrm{PI}\!\approx\!\mathrm{B2}$ as redundancy---%
$\mathbf{E}$ carrying what $\mathbf{P}$ carries---rather than as an encoder
failing to see what the pathologist named. It does not license the stronger
claim that a different encoder would leave every assignment intact:
representing the vocabulary and carrying the target's morphology are distinct
capacities, and a gene could move regime under an encoder that captures
target-relevant structure the vocabulary never named. That residual is the
honest limit of the reading, and the assignment is a property of the triple
(vocabulary, encoder, cohort) for exactly this reason.
Because the assignment turns on the \emph{sign} of a contrast rather than
on an absolute \auroc{}, the \emph{procedure} carries over to a new
cohort or modality with no per-cohort threshold to recalibrate, which
is what lets the same read-out run unchanged on glioma MRI and on the
electrocardiogram. That bears on the operationalisation and not on any
assignment it returns: the contrasts are computed from cohort-dependent
\auroc{}s, so a change of case mix can move a gene between regimes.
What we claim travels is the instrument, not its per-gene reading.

\paragraph{Why four conditions rather than a channel ablation.}
Reading a channel's contribution from the \auroc{} it costs when removed
has been criticised as a contribution measure, on the grounds that a
single drop conflates two different situations: a channel that is
informative in its own right, and a channel whose value exists only in
interaction with the other.
The four-condition design answers that objection by construction rather
than by argument. B2 and B3 expose each channel alone, \piabmil{} their
additive combination, and \fcmil{} their $2$-additive Choquet
combination, so the two situations separate: additive complementarity
raises \piabmil{} above both single channels while leaving
$\mathrm{FC}\!\approx\!\mathrm{PI}$, whereas interaction raises \fcmil{}
above \piabmil{} and is corroborated independently by the sign of the
learned $I_{jk}$, which is read off the fuzzy measure and not from any
\auroc{} at all. The single-channel ablation is accordingly used only
to corroborate the unique atoms it can legitimately evidence, and never
to adjudicate synergy.

The full mapping---every atom of \cref{eq:pid} against the model-side
signature that evidences it, the directional decision rule on the four
training conditions, the interpretability level that exposes it and the
regime it drives---is given in \cref{tab:atom_operationalization}. The signatures are
read, not numerically estimated; atom names follow the Bertschinger
convention.\cite{bertschinger2014unique}

\paragraph{Scope.}
No missing-channel correction\cite{choi2025icym2i} is needed: the
pattern channel $\mathbf{P}$ and the embedding channel $\mathbf{E}$ are
both computed from the same H\&E slide, so every case carries both. The
sampling-noise bound behind the directional policy is concrete: any
numerical atom would have to be built from the outputs of the two
single-channel models (B2, encoder only; B3, pattern only), whose
per-fold \auroc{} varies by $\pm 0.031$ to $\pm 0.081$ across folds
(\cref{tab:mutation_luad_resultssec}), and the differences to be
separated are no larger than that.

\paragraph{Relation to prior information-decomposition work.}
\Cref{eq:pid} is not new; what is new is the axis (two readings of one
modality rather than two modalities) and the purpose (adjudicating
whether a target is reachable at all, under declared power, rather than
choosing a fusion architecture). The position against the published
decomposition work in computational
pathology,\cite{liang2023quantifying,saisan2026ceilings} and the evidence
that direct atom estimators do not yet track downstream
performance,\cite{zhang2025evaluatingpid} are set out in Supplementary
Note~3.

% =====================================================================
\subsection*{Cross-domain instantiation protocols}
% =====================================================================

% Relocated from the Discussion.

The decomposition requires only an expert controlled vocabulary and an
unconstrained learned encoder of the same modality, both available far
from pathology. Cohorts, source ontologies and expected regime
assignments were fixed in advance (\cref{tab:crossdomain_protocol}).

\paragraph{Glioma MRI $\rightarrow$ genotype.}
VASARI is a consensus-defined controlled terminology of roughly thirty
semantic MR features---enhancement quality, necrosis, margin definition,
oedema, ependymal contact, grey-level heterogeneity---published with
per-subject scores and reported inter-observer agreement for the TCGA
glioma collections.\cite{gutman2013vasari,vasari_tcia} Of these, $26$
were retained after dropping features constant or unscored across the
cohort, and min--max rescaled. The encoder channel is a 2.5D CNN with
\abmil{} over axial slices of four co-registered sequences
(\cref{fig:glioma_example}), occupying the
position \ctranspath{} with tile-level \abmil{} holds in LUAD. IDH,
1p/19q and a patient-sex control are scored on the TCGA-LGG subjects
carrying real neuroradiologist VASARI scores, co-registered preoperative
MRI and a matched molecular call (up to $n\!=\!63$; IDH resolved on four
folds, one test fold single-class); that pilot is reported in
\cref{tab:glioma_domainb} and Supplementary Note~1 as directional reads
only. MGMT is scored on the full BraTS 2021
cohort\cite{baid2021brats} ($n\!=\!577$, ${\sim}52\%$ methylated), where
the VASARI channel is \emph{propagated} from the annotated seed rather
than expert-scored; the MGMT claim therefore rests on the unconstrained
encoder remaining at chance rather than on the strict empty-atom
separation that real VASARI licenses in LUAD. In both this arm and the
electrocardiogram the regime is adjudicated by the noise-aware rule
above, with the same seed-and-patient hierarchical bootstrap as LUAD:
each arm was trained with five seeds ($42$--$46$) and $2{,}000$ draws
resample seeds with replacement and then patients within the pooled
out-of-fold predictions of each drawn seed. The non-additive sub-form
(ii-b) of the 1p/19q complementarity is not resolved at $n\!\sim\!63$,
where the learned pair interactions shrink to
$\lvert I_{jk}\rvert\!\lesssim\!0.005$; we record that as deferred.

\paragraph{Electrocardiogram $\rightarrow$ physiological state.}
PTB-XL supplies $21{,}799$ records from $18{,}869$ patients annotated by
cardiologists with $71$ SCP-ECG statements, with the SNOMED-CT mapping
of PTB-XL$+$.\cite{wagner2020ptbxl,strodthoff2023ptbxlplus} Two
properties make this domain more favourable than histology: the source
vocabulary is already \emph{graded}, each statement carrying a
$0$--$100$ likelihood score extracted from the report, so the fuzzy
membership the pattern classifier must supply in LUAD is here given as
data; and the target set is an order of magnitude larger. The encoder channel is a compact 1D
convolutional network over the raw twelve-lead signal at $100$\,Hz
(\cref{fig:ecg_ptbxl}).
Patient sex and age are recorded for every case and are not derivable
from the statements, which is what makes them a leakage-free pair. Many
patients contribute more than one record, so the bootstrap resamples
patients and carries all of a patient's records together.
Whether the regimes remain distinguishable at $10^{2}$ targets rather
than six---as the published three-way partition of $108$ laboratory
indicators poses\cite{xiao2025anyecglab}---is not a question this run
settles.

% =====================================================================
\subsection*{Cohort-transfer protocol and controls (CPTAC-LUAD)}
% =====================================================================

% Relocated from the Discussion.

CPTAC-LUAD carries all six target genes and is a genuine distribution
shift rather than a resample of TCGA: it is enriched for never-smokers
and \textit{EGFR}-mutant, demographically more diverse
tumours.\cite{gillette2020cptacluad,cptacluad_data} We processed it
off-the-shelf---the frozen TCGA-trained pipeline applied to the $498$
slides from the $106$ genotyped patients with matched H\&E, with no stain
normalisation, colour calibration or domain adaptation, and slide scores
mean-pooled to the patient as stated in advance.

Three controls delimit the conclusion. First, aggregation: a slide-level
read-out ($n\!=\!498$), a max, a top-$2$/top-$3$ mean and a
single-largest-slide read-out all place \textit{TP53} between $0.49$ and
$0.57$, so the collapse is not an artefact of patient-level pooling.
Second, colour: re-running the transfer with Reinhard stain
normalisation of every tile to a representative TCGA-LUAD reference does
not recover it, so a global stain cast is not the driver and any
technical account must rest on deeper scanner- or site-level batch
signatures\cite{howard2021site} (Macenko normalisation remains untested).
Third, physical scale: the cohorts differ sharply in magnification
(TCGA-LUAD $76.4\%$ $40\times$, CPTAC-LUAD $98.6\%$ $20\times$), so
identical tiling at level~$0$ presents up to four times the tissue area
per tile; centre-cropping and rescaling every CPTAC tile to the field of
view of a $40\times$ TCGA tile lowers every condition and widens
fold-to-fold spread substantially, so scale mismatch is not the
mechanism, though the crop is itself lossy and the arm does not argue the
other way either.

The refit is within-cohort cross-validation and therefore establishes
that CPTAC slides carry the signal, not that one model serves both
cohorts. Regime adjudication follows the noise-aware rule above: the
deciding contrast $\piabmil{}-\mathrm{B2}$ is $+0.032$ with a
$2000$-sample paired-bootstrap $95\%$ CI of $-0.041$ to $+0.102$
natively, and $+0.073$ ($-0.038$ to $+0.174$) under field-of-view
matching. Both span zero, so the point-estimate read of (ii-a) is
withheld. Clearing the $n_+\!=\!32$ estimation floor says a model can be
\emph{fitted}; resolving a contrast between two fitted conditions is a
stricter demand, because what must exceed sampling error is the
\emph{difference} rather than either \auroc{}. Slide count does not
relieve it, since the genotype is a patient-level label and slides from
one tumour share it. Scaling the observed interval width to the observed
effect puts the cohort needed at roughly five times the present one, on
the order of $500$ genotyped patients.

% =====================================================================
\subsection*{Statistical analyses}
% =====================================================================

Five-fold patient-stratified cross-validation; folds were
created with mutation-stratified shuffling at the patient level to
preserve gene prevalence.
Per-fold metrics: \auroc{} (primary, prevalence-invariant),
\auprc{} (secondary, prevalence-sensitive),
positive class operating thresholds via Youden's
$J$.

\paragraph{Conclusive verdict criterion.}
The cohort-level Conclusive/Inconclusive flag
(\cref{tab:viz_folds_methods}) is defined as
\begin{equation}
\mathrm{Conclusive}(\mathrm{gene}) \iff
\frac{1}{5}\sum_{f=1}^{5}\auroc_{f}^{(\mathrm{gene})} \;\geq\; 0.70,
\label{eq:conclusive_verdict}
\end{equation}
where $\auroc_f^{(\mathrm{gene})}$ is the held-out AUROC of fold $f$
under the best MIL aggregator for that gene.
The cut is the conventional lower bound of \emph{acceptable discrimination}
for a fitted model,\cite{hosmer2000applied} and we treat it as a convention
rather than a derived quantity. It is therefore reported together with its
sensitivity: no LUAD target is Conclusive at $0.75$ or at $0.70$; lowering the
cut to $0.65$ flips \textit{TP53} on the patient-preserved protocol alone
($0.661$), and only at $0.60$ do \textit{TP53} and \textit{STK11} clear it
under both protocols (\cref{tab:threshold_sensitivity}). That no target in this cohort is
resolved is therefore a property of the estimates rather than of the
threshold, across the whole range a reader might reasonably prefer.
This is a population-level signal-presence flag, \emph{not} a
per-slide confidence threshold; the per-slide $P$-value (the
softmax-normalised mutation probability of the gene-specific head)
is reported independently and inherits the calibration of the
underlying \piabmil{} or \fcmil{} classifier.

\paragraph{Scope of the verdict.}
\Cref{eq:conclusive_verdict} thresholds a cohort mean, so the verdict is
a population statement: a Conclusive verdict licenses no confident
per-slide call, and an Inconclusive one does not imply that the modality
carries no information about the target, since the mean is taken over a
population that may be heterogeneous with respect to it. An Inconclusive
verdict is compatible with four distinct states---source-side atoms
empty, non-empty on a subpopulation, non-empty but unestimable, or a
heterogeneous gene-level label---with different remedies; the four
levels of analysis at which the quantities of this paper are estimated,
and the claims each supports, are set out in Supplementary Note~4 and
\cref{tab:levels_of_analysis}, and instantiated for the three modalities
in \cref{tab:levels_crossdomain}.

% =====================================================================
\subsection*{Released per-slide tooling}
% =====================================================================

The per-slide read-outs from which the model-side signatures above are
read are released with the code as a browser deployment
(\cref{fig:ui_stack}), numbered L1--L6 where the Supplementary tables
refer to them: the \abmil{} attention map (L1) and the
$\arg\max\mathbf{p}$ pattern overlay (L2) on the slide thumbnail; the
DeepSHAP\cite{lundberg2017shap} slide-level contribution split into its
embedding- and pattern-channel parts (L3); the per-slide Choquet--Shapley
values and pairwise interaction indices
$I_{jk}$ (L4);\cite{shapley1953value,grabisch1996application} the
single-channel ablation $\Delta P(\mathrm{mut})$ (L5); and a fixed
five-band trapezoidal mapping of each continuous score to a linguistic
category (L6), held constant across genes, slides and folds so that no
calibration data enter it. An optional narrative layer renders the bundle in prose through
a language-model agent restricted to those categorical labels and to
entities of a NCIt/MONDO knowledge
graph.\cite{sioutos2007ncit,mondo_shefchek2020} These read-outs are
computed at the slide and tile levels and are never aggregated into the
verdict (\cref{tab:levels_of_analysis}). The tooling is released as
software; its effect on a reader's decisions is not evaluated in this
study and no claim about it is made.

% =====================================================================
\subsection*{Reproducibility}
% =====================================================================

All experiments were run on a single H100 GPU (training) and
a T4 GPU (inference; the live \lchai{} deployment).
PyTorch 2.x, CTransPath checkpoint as published.\cite{wang2022ctranspath}
Seeds: 42, 142, 242, 342 and 442 for the five-seed mutation arm of
\cref{tab:mutation_luad_resultssec,tab:luad_contrasts}; the three-seed
control arm (\cref{tab:controls}) uses the first three. Each training
unit's seed is derived deterministically from the base seed, the
protocol, the gene, the fold and the condition, so every unit can be
regenerated from these five integers.
Mutation calls (MAF) were retrieved from the GDC on 03~April~2026.
Cohort, source-code and derived-artefact locations are given in the
Data and Code availability statements.

% Journal policy: use of an LLM must be documented in Methods.
\paragraph{Use of large language models.}
A language model is an optional component of the released per-slide
tooling (the narrative layer above, run with GPT-4 and Claude Sonnet);
none of the results reported here depends on it. In
preparing the manuscript language models were used for AI-assisted copy editing
only; no analysis, result, figure or reference was generated by a
language model, and no language model is an author.

% =====================================================================
\subsection*{Ethics statement}
% =====================================================================

The TCGA-LUAD cohort is a publicly available, de-identified
research dataset published under the controlled-access policy of
the NIH Genomic Data Commons.\cite{tcga2014luad}
The CPTAC-LUAD, TCGA-LGG/GBM, BraTS~2021 and PTB-XL cohorts of the
cohort-transfer and cross-domain arms are likewise public and
de-identified.
No new human-subject data were collected for this study, no
patient-identifiable information was accessed, and no further ethics
approval was required for this retrospective analysis of public data.

